\section{Discussion and Limitations}
\label{sec:discussion}
\subsection{Institutional value and responsibility}
An audit trail makes choices inspectable, but does not make the choices legitimate by itself. Constraints are set by actors with authority and incentives; versioning can expose changes but cannot prevent capture or selective compliance \cite{Selbst2019Abstraction,Metcalf2021AIA}. Institutions still need accessible appeal channels, a mandate for policy changes, appropriate staffing and a way to reverse or remedy decisions. A deployment should record who changed a policy, why, who may contest it and how resulting cases are handled.

Automation coverage measures automated outcomes, including rejection. It must not be equated with access to support or students helped. In screening mode the coverage floor can conflict structurally with limited K; increasing automation by relabeling cases changes the division of responsibility between software and staff. Our experiments expose that choice rather than endorse one taxonomy universally.

\subsection{What the evidence establishes}
The strongest evidence is computational: rules can be made explicit, decisions can be replayed on fixed scores, infeasibility can be surfaced, and revised exports can be gated. Sample-level constraints are not population guarantees. Policy selection on the same labelled audit set is exploratory. One historical model and split per dataset do not establish multi-seed predictive stability, comparative state of the art or educational treatment benefit.

The N=16 study is exploratory, uses author-defined standardized scenarios and mixed technical profiles, and has incomplete recruitment and training documentation. Volunteers and consent do not substitute for a formal ethics review, which was not obtained. No causal separation among provenance, structured information and replay is possible in its two conditions. Diagnosis and remediation correctness are perfectly concordant in these responses, so the study provides fewer independent outcome dimensions than its labels might suggest. A prospective study with independent cases, frontline stakeholders and documented institutional review is needed before making field claims.

\subsection{Uncertainty and data limitations}
Predictive entropy is a scalar confidence proxy, not a decomposition of epistemic and aleatoric uncertainty. Calibration and entropy/error correlation do not alone validate the safety of deferral. We therefore provide group-level selective-error summaries and explicitly conditional intervals. ASSISTments uses school-ID parity only to stress the constraint machinery; it does not measure protected-group disparities. OULAD uses binary gender and success labels, leaving wider intersectional and measurement questions unresolved \cite{BakerHawn2022AlgorithmicBiasEducation}.

\subsection{Privacy, integrity and operational boundaries}
Packets must contain the inputs required for the supported rule, with purpose-appropriate access and retention controls. Publishing medical information or raw feature histories is not required by our minimum contract. An institution can keep restricted source evidence behind resolvable references, but missing evidence must be reported rather than treated as replay success. Hashes establish consistency of bytes, not external authenticity or truth. The revised runner's review-required file is not a connected case-management process.

We make no automatic legal-compliance claim. A particular deployment requires jurisdiction-specific institutional review. Instrumented per-case audit cost, staffing capacity, actual access to appeal, independent prospective evaluation and model re-scoring under corrected inputs remain future work. Changing a feature or target creates a new inference or modeling task and must not be described as a policy-only replay.
